\subsection{Comments and generalizations} 

\subsubsection{Extended symmetry} The $W$-action on $\cohog{*}{\sB_{e}}$ can be extended to an action of a larger algebra in various ways, if we use more sophisticated cohomology theories. On the equivariant cohomology $\upH^{*}_{\tilG_{e}}(\sB_{e})$, there is an action of the {\em graded affine Hecke algebra} (see Lusztig \cite{LuGrHk}). On the $\tilG_{e}$-equivariant $K$-group of $\sB_{e}$, there is an action of the {\em affine Hecke algebra} (see Kazhdan-Lusztig \cite{KLAffHk} and Chriss-Ginzburg \cite{CG}). 

\subsubsection{The group version} There are obvious analogs of the Springer resolution and the Grothendieck alteration when $\sN$ and $\frg$ are replaced with the unipotent variety $\sU\subset G$ and $G$ itself. When $\chark$ is large, the exponential map identifies $\sN$ with $\sU$ in a $G$-equivariant manner, hence the theories of Springer fibers for nilpotent elements and unipotent elements are identical. The group version of the perverse sheaf $\calS_{\frg}$ and its irreducible direct summands are precursors of {\em character sheaves}, a theory developed by Lusztig (\cite{LuChShI}, \cite{LuChShII}, \cite{LuChShIV} and \cite{LuChShV}) to study characters of the finite groups  $G(\FF_{q})$.

\subsubsection{Partial Springer resolutions}\label{sss:partial Spr} We may define analogs of $\sB_{e}$ in partial flag varieties.  Let $\sP$ be a partial flag variety of $G$ classifying parabolic subgroups $P$ of $G$ of a given type. There are two analogs of the map $\pi:\tN\to\sN$ one may consider. 

First,  instead of considering $\tN=T^{*}\sB$, we may consider $T^{*}\sP$, which classifies pairs $(e, P)\in\sN\times\sP$ such that $e\in \Lie\frn_{P}$, where $\frn_{P}$ is the nilpotent radical of $\Lie P$. Let $\tau_{\sP}: T^{*}\sP\to\sN$ be the first projection. In general this map is not surjective, its image is the closure of a nilpotent orbit $\calO_{\calP}$. The orbit $\calO_{P}$ is characterized by the property that its intersection with $\frn_{P}$ is dense in $\frn_{P}$, for any $P\in\sP$. This is called the {\em Richardson class} associated to parabolic subgroups of type $\sP$. When $G$ is of type $A$, each nilpotent class $\calO$ is the Richardson class associated to parabolic subgroups of some type $\sP$ (not unique in general). The map $T^{*}\sP\to \ov\calO$ is a resolution of singularities. For general $G$, not every nilpotent orbit is Richardson.

Second, we may consider the subscheme $\tN_{\sP}\subset\sN\times \sP$ classifying pairs $(e,P)$ such that $e\in\sN_{P}$, where $\sN_{P}\subset\Lie P$ is the nilpotent cone of $P$.  The projection $\pi_{\sP}: \tN_{\sP}\to \sN$ is now surjective, and is a partial resolution of singularities. The Springer resolution $\pi$ can be factored as
\begin{equation*}
\pi: \tN=\tN_{\sB}\xrightarrow{\nu_{\sP}}\tN_{\sP}\xrightarrow{\pi_{\sP}} \sN. 
\end{equation*}

We have an embedding $T^{*}\sP\incl \tN_{\sP}$.  We may consider the fibers of either $\tau_{\sP}$ or $\pi_{\sP}$ as parabolic analogs of Springer fibers.  We call them {\em partial Springer fibers}. The Springer action of $W$ on the cohomology of $\sB_{e}$ has an analog for partial Springer fibers. For more information, we refer the readers to \cite{BM}.

\subsubsection{Hessenberg varieties} The Grothendieck alteration $\pi_{\frg}:\tg\to \frg$ admits a generalization where $\frg$ is replaced with another linear representation of $G$.

Fix a Borel subgroup $B$ of $G$. Let $(V,\rho)$ be a representation of $G$ and $V^{+}\subset V$ be a subspace which is stable under $B$. Now we use the pair $(V,V^{+})$ instead of the pair $(\frg,\frb)$, we get a generalization of the Grothendieck alteration. More precisely, let $\tilV\subset V\times \calB$ be the subscheme consisting of pairs $(v, gB)\in V\times\sB$ such that $v\in \rho(g)V^{+}$. Let $\pi_{V}:\tilV\to V$ be the first projection. The fibers of $\pi_{V}$ are called {\em Hessenberg varieties}.

Hessenberg varieties appear naturally in the study of certain affine Springer fibers, as we will see in \S\ref{ss:aff examples}. For more information on Hessenberg varieties, see \cite{GKM} and \cite{OY}.


\subsection{Exercises}

\subsubsection{} For $G=\SL_{n}$, determine the sizes of the Jordan blocks of a regular and subregular nilpotent element of $\frg$.
 
\subsubsection{}\label{exI:nonred Spr} For $G=\SL_{2}$ and $\SL_{3}$, calculate the coordinate ring of the {\em non-reduced} Springer fiber $\tB_{e}$ for a regular nilpotent element $e$. Show also that the Springer fiber $\sB_{e}$ is indeed a reduced point. 

{\em Hint}: if you write $e$ as an upper triangular matrix, then $\tB_{e}$ lies in the big Bruhat cell of the flag variety $\sB$, from which you get coordinates for your calculation.


\subsubsection{}\label{exI:211Sp4} Verify the statement in \S\ref{sss:Sp4}: consider $G=\Sp(V)$, $V=\jiao{v_{1},v_{2},v_{3},v_{4}}$ with the symplectic form given by $\jiao{v_{i},v_{5-i}}=1$ if $i=1,2$ and $\jiao{v_{i},v_{j}}=0$ for $i+j\neq 5$. Let $e: v_{4}\mapsto v_{1}\to 0, v_{2}\mapsto 0,v_{3}\mapsto 0$. Then any flag $0\subset V_{1}\subset V_{2}\subset V_{1}^{\bot}\subset V$ in $\sB_{e}$ must satisfy
\begin{equation*}
\jiao{v_{1}}\subset V_{2}\subset\jiao{v_{1},v_{2},v_{3}}.
\end{equation*}
Moreover, this is the only condition for a flag $0\subset V_{1}\subset V_{2}\subset V_{1}^{\bot}\subset V$ to lie in $\sB_{e}$.


\subsubsection{}\label{exI:inv} Let $e\in\sN$.  Let $B\subset G$ be a Borel subgroup.
\begin{enumerate}
\item  Let $\alpha$ be a simple root. Let $P_{\alpha}\supset B$ be a parabolic subgroup whose Levi factor has semisimple rank one with roots $\pm\alpha$. Let $\sP_{\alpha}\cong G/P_{\alpha}$ be the partial flag variety of $G$ classifying parabolics conjugate to $P_{\alpha}$. Restricting the projection $\sB\to \sP_{\alpha}$  to $\sB_{e}$, we get a map
\begin{equation*}
\pi_{\alpha}: \sB_{e}\to \pi_{\alpha}(\sB_{e}).
\end{equation*}
Show that the pullback $\pi^{*}_{\alpha}$ on cohomology induces an isomorphism
\begin{equation}\label{sa inv}
\cohog{*}{\pi_{\alpha}(\sB_{e})}\cong \cohog{*}{\sB_{e}}^{s_{\alpha}}
\end{equation}
where $s_{\alpha}\in W$ is the simple reflection associated with $\alpha$, which acts on $\cohog{*}{\sB_{e}}$ via Springer's action.
\item Can you generalize \eqref{sa inv} to other partial flag varieties?
\end{enumerate}


\subsubsection{}\label{exI:SL3}  Let $G=\SL_{3}$ and let $e\in\sN$ be a subregular element. Calculate the action of $S_{3}$ on the two dimensional $\cohog{2}{\sB_{e}}$ in terms of the basis given by the fundamental classes of the two irreducible components (see \S\ref{sss:SL3 subreg}).

{\em Hint}: for this and the next problem, you may find Exercise \ref{exI:inv} useful.

\subsubsection{}\label{exI:Sp4} Describe the Springer fibers for $\Sp_{4}$. Calculate the Springer correspondence for $G=\Sp_{4}$ explicitly. 


\subsubsection{} Using the dimension formula for $\sB_{e}$, verify that the Springer resolution $\pi$ is semismall.

\subsubsection{}\label{exI:dim BX} Let $X\in\frg$ and let $X=X_{s}+X_{n}$ be the Jordan decomposition of $X$. Let $H=C_{G}(X_{s})\subset G$. This is a Levi subgroup of $G$. Let $\sB^{H}$ be the flag variety of $H$ and let $\sB^{H}_{X_{n}}$ be the Springer fiber of $X_{n}$ viewed as a nilpotent element in $\Lie H$. Show that $\dim\tB_{X}=\dim \sB^{H}_{X_{n}}$.


\subsubsection{} Use Exercise \ref{exI:dim BX} and the dimension formula \eqref{Spr dim} to derive a formula for the dimension of $\tB_{X}$ for all elements $X\in\frg$. Use your formula to prove that the Grothendieck alteration $\pi_{\frg}$ is small.


\subsubsection{} Show that $\sN$ is rationally smooth; i.e., its intersection cohomology complex is isomorphic to the shifted constant sheaf $\Ql[\dim\sN]$.

{\em Hint}: the largest direct summand in $\bR\pi_{*}\Ql$ is the shifted IC sheaf of $\sN$, and it is also the restriction of a direct summand of $\bR\pi_{\frg,*}\Ql$.


\subsubsection{}\label{exI:conn} Show that the Springer fibers $\sB_{e}$ are connected. 

{\em Hint}: the $\upH^{0}$ of the Springer fibers are packed in some sheaf.


\subsubsection{} Denote the simple roots of $G$ by $\{\alpha_{1},\cdots,\alpha_{r}\}$. A parabolic subgroup $P\subset G$ is called of type $i$ if the roots of its Levi quotient $L_{P}$ are $\pm\alpha_{i}$.
\begin{enumerate}
\item Let $1\leq i\leq r$ and let $P$ be a parabolic subgroup of type $i$. Let $\frn_{P}$ be the nilpotent radical of $\Lie P$. Show that $\frn_{P}\cap \calO_{\subreg}$ is dense in $\frn_{P}$.
\item Let $e\in\calO_{\subreg}$. Show that for each $i$, there are finitely many parabolics $P$ of type $i$ such that $e\in\frn_{P}$. For each such $P$, the subvariety $C_{P}:=\{B\in\sB_{e}|B\subset P\}$ of $\sB_{e}$ is isomorphic to $\PP^{1}$, and is called a curve of type $i$.
\item For parabolics $P\neq Q$ of type $i$ and $j$, show that $C_{P}\cap C_{Q}$ is either empty or a point.
\item Let $P$ be a parabolic subgroup of type $i$ such that $e\in\frn_{P}$. For any $1\leq j\leq r, i\neq j$, $C_{P}$ intersects exactly $-\jiao{\alpha^{\vee}_{i}, \alpha_{j}}$ curves of type $j$.
\item Show that $\sB_{e}$ has the configuration described in \S\ref{sss:subreg}.
\item Use Exercise \ref{exI:inv} to calculate the Springer action of $W$ on $\cohog{2}{\sB_{e}}$.
\end{enumerate}



%%%%%%%%%%%%%%%%%%%%%%%



